\newcommand{\coursemodule}{Module 06}
% ════════════════════════════════════════════════════════════════════════════
% Pixels, Frames & Transformers: shared Beamer preamble (WUR style)
% Adapted from the preamble of Bags, Vectors & Transformers.
%
% Use in a module deck (from <part>/<module>/slides/):
%   \newcommand{\coursemodule}{Module 01}
%   % ════════════════════════════════════════════════════════════════════════════
% Pixels, Frames & Transformers: shared Beamer preamble (WUR style)
% Adapted from the preamble of Bags, Vectors & Transformers.
%
% Use in a module deck (from <part>/<module>/slides/):
%   \newcommand{\coursemodule}{Module 01}
%   % ════════════════════════════════════════════════════════════════════════════
% Pixels, Frames & Transformers: shared Beamer preamble (WUR style)
% Adapted from the preamble of Bags, Vectors & Transformers.
%
% Use in a module deck (from <part>/<module>/slides/):
%   \newcommand{\coursemodule}{Module 01}
%   \input{../../../shared/preamble.tex}
%   \begin{document}
%   \moduletitle{Images as Data}
%   ...
% Compile with: latexmk -pdf <deck>.tex
% ════════════════════════════════════════════════════════════════════════════

\documentclass[aspectratio=169]{beamer}

% ── Packages ────────────────────────────────────────────────────────────────
\usepackage[T1]{fontenc}
\usepackage[utf8]{inputenc}
\usepackage{lmodern}
\usepackage{microtype}
\usepackage{xcolor}
\usepackage{tikz}
\usetikzlibrary{arrows.meta, positioning, calc, shapes.geometric, decorations.pathreplacing, matrix, fit}
\usepackage{booktabs}
\usepackage{fontawesome5}
\usepackage{array}
\usepackage[most]{tcolorbox}

\providecommand{\coursemodule}{}

% ── WUR color palette ───────────────────────────────────────────────────────
\definecolor{wurgreen}{HTML}{34B233}
\definecolor{darkgreen}{HTML}{1F6B1E}
\definecolor{lightgreen}{HTML}{E8F7E8}
\definecolor{darknavy}{HTML}{1A1A2E}
\definecolor{midgray}{HTML}{6B7280}
\definecolor{lightgray}{HTML}{F3F4F6}
\definecolor{bodytext}{HTML}{1F2937}
\definecolor{codegreen}{HTML}{9AE6B4}
\definecolor{warnred}{HTML}{FEE2E2}
\definecolor{warnborder}{HTML}{EF4444}

% ── Beamer theme ────────────────────────────────────────────────────────────
\usetheme{default}
\usecolortheme{default}
\usefonttheme{professionalfonts}
\setbeamertemplate{navigation symbols}{}
\setbeamertemplate{itemize items}{\textcolor{wurgreen}{\textbullet}}
\setbeamertemplate{itemize subitem}{\textcolor{midgray}{--}}
\setbeamertemplate{enumerate items}{\textcolor{wurgreen}{\textbf{\insertenumlabel.}}}
\setbeamercolor{frametitle}{fg=white, bg=darknavy}
\setbeamertemplate{frametitle}{%
  \begin{beamercolorbox}[wd=\paperwidth, ht=1.1cm, dp=0.2cm, leftskip=0.6cm]{frametitle}
    \usebeamerfont{frametitle}\insertframetitle
  \end{beamercolorbox}%
}
\setbeamerfont{frametitle}{size=\large, series=\bfseries}
\setbeamercolor{background canvas}{bg=white}
\setbeamercolor{normal text}{fg=bodytext}
\setbeamerfont{normal text}{size=\small}
\setbeamertemplate{footline}{%
  \begin{beamercolorbox}[wd=\paperwidth, ht=0.4cm, dp=0.15cm]{footline}%
    \hspace{0.5cm}%
    \textcolor{midgray}{\tiny Pixels, Frames \& Transformers · \coursemodule\ · Wageningen University \& Research}%
    \hfill%
    \textcolor{midgray}{\tiny \insertframenumber}%
    \hspace{0.4cm}%
  \end{beamercolorbox}%
}

% ── Boxes ───────────────────────────────────────────────────────────────────
\tcbset{wurbase/.style={arc=3pt, boxrule=0.6pt, left=7pt, right=7pt, top=5pt, bottom=5pt,
                        fontupper=\small, before skip=4pt, after skip=4pt}}
\newtcolorbox{greenbox}[1][]{wurbase, colback=lightgreen, colframe=wurgreen, #1}
\newtcolorbox{graybox}[1][]{wurbase, colback=lightgray, colframe=lightgray, #1}
\newtcolorbox{navybox}[1][]{wurbase, colback=darknavy, colframe=darknavy, colupper=white, #1}
\newtcolorbox{warnbox}[1][]{wurbase, colback=warnred, colframe=warnborder, #1}
\newtcolorbox{codebox}[1][]{wurbase, colback=darknavy, colframe=darknavy, colupper=white,
                            fontupper=\ttfamily\scriptsize, left=6pt, top=4pt, bottom=4pt, #1}

% ── Title slide ─────────────────────────────────────────────────────────────
\newcommand{\moduletitle}[1]{%
{
\setbeamertemplate{footline}{}
\begin{frame}[plain]
  \begin{tikzpicture}[remember picture, overlay]
    \fill[darknavy] (current page.south west) rectangle (current page.north east);
    \fill[wurgreen] (current page.south west) rectangle ([xshift=0.45cm]current page.north west);
    \node[anchor=west, text=white, font=\huge\bfseries, text width=13cm]
      at ([xshift=1.1cm, yshift=1.2cm]current page.west) {Pixels, Frames \& Transformers};
    \node[anchor=west, text=wurgreen, font=\large, text width=13cm]
      at ([xshift=1.1cm, yshift=0.05cm]current page.west) {#1 · \coursemodule};
    \node[anchor=west, text=midgray, font=\small, text width=12cm]
      at ([xshift=1.1cm, yshift=-1.5cm]current page.west)
      {Denise J. Roth, PhD Candidate\\
       Strategic Communication Group · Wageningen University \& Research};
  \end{tikzpicture}
\end{frame}
}}

% ── Small helpers ───────────────────────────────────────────────────────────
\newcommand{\keyterm}[1]{\textbf{\textcolor{darkgreen}{#1}}}
\newcommand{\muted}[1]{\textcolor{midgray}{#1}}
\newcommand{\credit}[1]{{\tiny\textcolor{midgray}{#1}}}
\newcommand{\notebookslide}[1]{%
\begin{frame}{In the Notebook}
\vspace{0.1cm}
#1
\end{frame}}
\newcommand{\nbitem}[2]{\item[\textcolor{wurgreen}{#1}] #2}

%   \begin{document}
%   \moduletitle{Images as Data}
%   ...
% Compile with: latexmk -pdf <deck>.tex
% ════════════════════════════════════════════════════════════════════════════

\documentclass[aspectratio=169]{beamer}

% ── Packages ────────────────────────────────────────────────────────────────
\usepackage[T1]{fontenc}
\usepackage[utf8]{inputenc}
\usepackage{lmodern}
\usepackage{microtype}
\usepackage{xcolor}
\usepackage{tikz}
\usetikzlibrary{arrows.meta, positioning, calc, shapes.geometric, decorations.pathreplacing, matrix, fit}
\usepackage{booktabs}
\usepackage{fontawesome5}
\usepackage{array}
\usepackage[most]{tcolorbox}

\providecommand{\coursemodule}{}

% ── WUR color palette ───────────────────────────────────────────────────────
\definecolor{wurgreen}{HTML}{34B233}
\definecolor{darkgreen}{HTML}{1F6B1E}
\definecolor{lightgreen}{HTML}{E8F7E8}
\definecolor{darknavy}{HTML}{1A1A2E}
\definecolor{midgray}{HTML}{6B7280}
\definecolor{lightgray}{HTML}{F3F4F6}
\definecolor{bodytext}{HTML}{1F2937}
\definecolor{codegreen}{HTML}{9AE6B4}
\definecolor{warnred}{HTML}{FEE2E2}
\definecolor{warnborder}{HTML}{EF4444}

% ── Beamer theme ────────────────────────────────────────────────────────────
\usetheme{default}
\usecolortheme{default}
\usefonttheme{professionalfonts}
\setbeamertemplate{navigation symbols}{}
\setbeamertemplate{itemize items}{\textcolor{wurgreen}{\textbullet}}
\setbeamertemplate{itemize subitem}{\textcolor{midgray}{--}}
\setbeamertemplate{enumerate items}{\textcolor{wurgreen}{\textbf{\insertenumlabel.}}}
\setbeamercolor{frametitle}{fg=white, bg=darknavy}
\setbeamertemplate{frametitle}{%
  \begin{beamercolorbox}[wd=\paperwidth, ht=1.1cm, dp=0.2cm, leftskip=0.6cm]{frametitle}
    \usebeamerfont{frametitle}\insertframetitle
  \end{beamercolorbox}%
}
\setbeamerfont{frametitle}{size=\large, series=\bfseries}
\setbeamercolor{background canvas}{bg=white}
\setbeamercolor{normal text}{fg=bodytext}
\setbeamerfont{normal text}{size=\small}
\setbeamertemplate{footline}{%
  \begin{beamercolorbox}[wd=\paperwidth, ht=0.4cm, dp=0.15cm]{footline}%
    \hspace{0.5cm}%
    \textcolor{midgray}{\tiny Pixels, Frames \& Transformers · \coursemodule\ · Wageningen University \& Research}%
    \hfill%
    \textcolor{midgray}{\tiny \insertframenumber}%
    \hspace{0.4cm}%
  \end{beamercolorbox}%
}

% ── Boxes ───────────────────────────────────────────────────────────────────
\tcbset{wurbase/.style={arc=3pt, boxrule=0.6pt, left=7pt, right=7pt, top=5pt, bottom=5pt,
                        fontupper=\small, before skip=4pt, after skip=4pt}}
\newtcolorbox{greenbox}[1][]{wurbase, colback=lightgreen, colframe=wurgreen, #1}
\newtcolorbox{graybox}[1][]{wurbase, colback=lightgray, colframe=lightgray, #1}
\newtcolorbox{navybox}[1][]{wurbase, colback=darknavy, colframe=darknavy, colupper=white, #1}
\newtcolorbox{warnbox}[1][]{wurbase, colback=warnred, colframe=warnborder, #1}
\newtcolorbox{codebox}[1][]{wurbase, colback=darknavy, colframe=darknavy, colupper=white,
                            fontupper=\ttfamily\scriptsize, left=6pt, top=4pt, bottom=4pt, #1}

% ── Title slide ─────────────────────────────────────────────────────────────
\newcommand{\moduletitle}[1]{%
{
\setbeamertemplate{footline}{}
\begin{frame}[plain]
  \begin{tikzpicture}[remember picture, overlay]
    \fill[darknavy] (current page.south west) rectangle (current page.north east);
    \fill[wurgreen] (current page.south west) rectangle ([xshift=0.45cm]current page.north west);
    \node[anchor=west, text=white, font=\huge\bfseries, text width=13cm]
      at ([xshift=1.1cm, yshift=1.2cm]current page.west) {Pixels, Frames \& Transformers};
    \node[anchor=west, text=wurgreen, font=\large, text width=13cm]
      at ([xshift=1.1cm, yshift=0.05cm]current page.west) {#1 · \coursemodule};
    \node[anchor=west, text=midgray, font=\small, text width=12cm]
      at ([xshift=1.1cm, yshift=-1.5cm]current page.west)
      {Denise J. Roth, PhD Candidate\\
       Strategic Communication Group · Wageningen University \& Research};
  \end{tikzpicture}
\end{frame}
}}

% ── Small helpers ───────────────────────────────────────────────────────────
\newcommand{\keyterm}[1]{\textbf{\textcolor{darkgreen}{#1}}}
\newcommand{\muted}[1]{\textcolor{midgray}{#1}}
\newcommand{\credit}[1]{{\tiny\textcolor{midgray}{#1}}}
\newcommand{\notebookslide}[1]{%
\begin{frame}{In the Notebook}
\vspace{0.1cm}
#1
\end{frame}}
\newcommand{\nbitem}[2]{\item[\textcolor{wurgreen}{#1}] #2}

%   \begin{document}
%   \moduletitle{Images as Data}
%   ...
% Compile with: latexmk -pdf <deck>.tex
% ════════════════════════════════════════════════════════════════════════════

\documentclass[aspectratio=169]{beamer}

% ── Packages ────────────────────────────────────────────────────────────────
\usepackage[T1]{fontenc}
\usepackage[utf8]{inputenc}
\usepackage{lmodern}
\usepackage{microtype}
\usepackage{xcolor}
\usepackage{tikz}
\usetikzlibrary{arrows.meta, positioning, calc, shapes.geometric, decorations.pathreplacing, matrix, fit}
\usepackage{booktabs}
\usepackage{fontawesome5}
\usepackage{array}
\usepackage[most]{tcolorbox}

\providecommand{\coursemodule}{}

% ── WUR color palette ───────────────────────────────────────────────────────
\definecolor{wurgreen}{HTML}{34B233}
\definecolor{darkgreen}{HTML}{1F6B1E}
\definecolor{lightgreen}{HTML}{E8F7E8}
\definecolor{darknavy}{HTML}{1A1A2E}
\definecolor{midgray}{HTML}{6B7280}
\definecolor{lightgray}{HTML}{F3F4F6}
\definecolor{bodytext}{HTML}{1F2937}
\definecolor{codegreen}{HTML}{9AE6B4}
\definecolor{warnred}{HTML}{FEE2E2}
\definecolor{warnborder}{HTML}{EF4444}

% ── Beamer theme ────────────────────────────────────────────────────────────
\usetheme{default}
\usecolortheme{default}
\usefonttheme{professionalfonts}
\setbeamertemplate{navigation symbols}{}
\setbeamertemplate{itemize items}{\textcolor{wurgreen}{\textbullet}}
\setbeamertemplate{itemize subitem}{\textcolor{midgray}{--}}
\setbeamertemplate{enumerate items}{\textcolor{wurgreen}{\textbf{\insertenumlabel.}}}
\setbeamercolor{frametitle}{fg=white, bg=darknavy}
\setbeamertemplate{frametitle}{%
  \begin{beamercolorbox}[wd=\paperwidth, ht=1.1cm, dp=0.2cm, leftskip=0.6cm]{frametitle}
    \usebeamerfont{frametitle}\insertframetitle
  \end{beamercolorbox}%
}
\setbeamerfont{frametitle}{size=\large, series=\bfseries}
\setbeamercolor{background canvas}{bg=white}
\setbeamercolor{normal text}{fg=bodytext}
\setbeamerfont{normal text}{size=\small}
\setbeamertemplate{footline}{%
  \begin{beamercolorbox}[wd=\paperwidth, ht=0.4cm, dp=0.15cm]{footline}%
    \hspace{0.5cm}%
    \textcolor{midgray}{\tiny Pixels, Frames \& Transformers · \coursemodule\ · Wageningen University \& Research}%
    \hfill%
    \textcolor{midgray}{\tiny \insertframenumber}%
    \hspace{0.4cm}%
  \end{beamercolorbox}%
}

% ── Boxes ───────────────────────────────────────────────────────────────────
\tcbset{wurbase/.style={arc=3pt, boxrule=0.6pt, left=7pt, right=7pt, top=5pt, bottom=5pt,
                        fontupper=\small, before skip=4pt, after skip=4pt}}
\newtcolorbox{greenbox}[1][]{wurbase, colback=lightgreen, colframe=wurgreen, #1}
\newtcolorbox{graybox}[1][]{wurbase, colback=lightgray, colframe=lightgray, #1}
\newtcolorbox{navybox}[1][]{wurbase, colback=darknavy, colframe=darknavy, colupper=white, #1}
\newtcolorbox{warnbox}[1][]{wurbase, colback=warnred, colframe=warnborder, #1}
\newtcolorbox{codebox}[1][]{wurbase, colback=darknavy, colframe=darknavy, colupper=white,
                            fontupper=\ttfamily\scriptsize, left=6pt, top=4pt, bottom=4pt, #1}

% ── Title slide ─────────────────────────────────────────────────────────────
\newcommand{\moduletitle}[1]{%
{
\setbeamertemplate{footline}{}
\begin{frame}[plain]
  \begin{tikzpicture}[remember picture, overlay]
    \fill[darknavy] (current page.south west) rectangle (current page.north east);
    \fill[wurgreen] (current page.south west) rectangle ([xshift=0.45cm]current page.north west);
    \node[anchor=west, text=white, font=\huge\bfseries, text width=13cm]
      at ([xshift=1.1cm, yshift=1.2cm]current page.west) {Pixels, Frames \& Transformers};
    \node[anchor=west, text=wurgreen, font=\large, text width=13cm]
      at ([xshift=1.1cm, yshift=0.05cm]current page.west) {#1 · \coursemodule};
    \node[anchor=west, text=midgray, font=\small, text width=12cm]
      at ([xshift=1.1cm, yshift=-1.5cm]current page.west)
      {Denise J. Roth, PhD Candidate\\
       Strategic Communication Group · Wageningen University \& Research};
  \end{tikzpicture}
\end{frame}
}}

% ── Small helpers ───────────────────────────────────────────────────────────
\newcommand{\keyterm}[1]{\textbf{\textcolor{darkgreen}{#1}}}
\newcommand{\muted}[1]{\textcolor{midgray}{#1}}
\newcommand{\credit}[1]{{\tiny\textcolor{midgray}{#1}}}
\newcommand{\notebookslide}[1]{%
\begin{frame}{In the Notebook}
\vspace{0.1cm}
#1
\end{frame}}
\newcommand{\nbitem}[2]{\item[\textcolor{wurgreen}{#1}] #2}


\begin{document}

\moduletitle{Transfer Learning and Image Embeddings}

% ─────────────────────────────────────────────────────────────────────────────
\begin{frame}{Don't Start from Scratch}
A network trained on millions of images has already learned edges, textures, parts and objects.
\keyterm{Transfer learning} reuses that knowledge, just as BERT is reused for text.

\vspace{0.25cm}
\begin{columns}[T]
\begin{column}{0.48\textwidth}
\begin{greenbox}
\textbf{1. Off the shelf}\\
Use the model's own labels: an ImageNet model names 1,000 categories, from \emph{tabby cat} to
\emph{space shuttle}.
\end{greenbox}
\end{column}
\begin{column}{0.48\textwidth}
\begin{greenbox}
\textbf{2. As an embedding machine}\\
Remove the last layer, and the model turns every image into a vector. Train your own classifier
on top, with your categories.
\end{greenbox}
\end{column}
\end{columns}

\vspace{0.2cm}
\begin{graybox}
\textbf{ImageNet}: about 1.3 million training photos in 1,000 categories, the benchmark behind a
decade of computer vision. Its categories were chosen by computer scientists: 120 of them are dog
breeds, and none is \emph{protest}, \emph{politician} or \emph{crowd}.
\end{graybox}
\end{frame}

% ─────────────────────────────────────────────────────────────────────────────
\begin{frame}{Cut Off the Last Layer: Image Embeddings}
\vspace{0.2cm}
\centering
\begin{tikzpicture}[
  img/.style={draw=midgray, fill=lightgray, rounded corners=3pt, font=\footnotesize, align=center,
              inner sep=5pt, text width=1.4cm, minimum height=1.2cm},
  net/.style={draw=wurgreen, fill=lightgreen, rounded corners=3pt, font=\footnotesize, align=center,
              inner sep=5pt, text width=2.9cm, minimum height=1.2cm},
  emb/.style={draw=darknavy, fill=darknavy, text=white, rounded corners=3pt, font=\footnotesize\bfseries,
              align=center, inner sep=5pt, text width=2.2cm, minimum height=1.2cm},
  cut/.style={draw=warnborder, fill=warnred, rounded corners=3pt, font=\footnotesize, align=center,
              inner sep=5pt, text width=3.0cm, minimum height=1.2cm, dashed},
  arr/.style={-{Stealth[length=2mm]}, thick, color=midgray}]
\node[img] (i) {an image};
\node[net, right=0.5cm of i] (n) {pre-trained layers\\ \muted{frozen}};
\node[emb, right=0.5cm of n] (e) {embedding\\ \normalfont 2,048 numbers};
\node[cut, right=0.5cm of e] (c) {1,000 ImageNet\\ scores: \textbf{removed}};
\draw[arr] (i) -- (n); \draw[arr] (n) -- (e); \draw[arr, dashed] (e) -- (c);
\end{tikzpicture}

\vspace{0.35cm}
\raggedright
\begin{columns}[T]
\begin{column}{0.48\textwidth}
\begin{graybox}
\textbf{ResNet-50}: a CNN trained \emph{with labels} on ImageNet. 2,048 numbers per image.
\end{graybox}
\end{column}
\begin{column}{0.48\textwidth}
\begin{greenbox}
\textbf{DINOv2} \muted{(Oquab et al., 2024)}: a vision transformer trained \emph{without labels}
on 142 million images. 384 numbers per image (small version).
\end{greenbox}
\end{column}
\end{columns}
\muted{Images that look alike, to the model, get vectors that point in similar directions. With
vectors of length 1, their dot product is the \textbf{cosine similarity}.}
\end{frame}

% ─────────────────────────────────────────────────────────────────────────────
\begin{frame}{Three Ways to Use a Pre-trained Model}
\vspace{0.1cm}
{\footnotesize
\begin{tabular}{@{}>{\raggedright\arraybackslash}p{3.0cm}>{\raggedright\arraybackslash}p{3.5cm}>{\raggedright\arraybackslash}p{3.0cm}>{\raggedright\arraybackslash}p{3.0cm}@{}}
\toprule
& \textbf{What is trained} & \textbf{Cost} & \textbf{With few labels} \\
\midrule
\textbf{Linear probe} & a logistic regression on the frozen embeddings & seconds, no GPU once embedded & the safe default \\[3pt]
\textbf{Fine-tune the last block} & the last block plus a new output layer & minutes on a GPU; more settings & sometimes better, often not \\[3pt]
\textbf{Fine-tune everything} & every weight, at a low learning rate & the most compute and care & overfits easily \\
\bottomrule
\end{tabular}}

\vspace{0.35cm}
\begin{greenbox}
\textbf{Start with a linear probe on good embeddings.} It is the image version of a classifier on
sentence embeddings, and usually far better than a CNN trained from scratch on the same few hundred
labels. Fine-tune only when you have more labels and the probe has clearly hit a ceiling.
\end{greenbox}
\end{frame}

% ─────────────────────────────────────────────────────────────────────────────
\begin{frame}{Choose Settings Without Touching the Test Set}
\begin{columns}[T]
\begin{column}{0.5\textwidth}
\centering
\begin{tikzpicture}[
  tr/.style={draw=wurgreen, fill=lightgreen, rounded corners=3pt, font=\footnotesize, align=center,
             inner sep=5pt, text width=5.2cm},
  te/.style={draw=darknavy, fill=darknavy, text=white, rounded corners=3pt, font=\footnotesize,
             align=center, inner sep=5pt, text width=5.2cm},
  arr/.style={-{Stealth[length=2mm]}, thick, color=midgray}]
\node[tr] (a) {\textbf{training images}\\ try $C$ = 0.01, 0.1, 1, 10, 100\\ with 5-fold cross-validation};
\node[tr, below=0.35cm of a] (b) {refit with the best setting};
\node[te, below=0.35cm of b] (c) {\textbf{your 40 coded images}\\ one final score};
\draw[arr] (a) -- (b); \draw[arr] (b) -- (c);
\end{tikzpicture}
\end{column}
\begin{column}{0.47\textwidth}
Picking a setting by its score on the test set turns the test set into training data: you pick
whatever happened to fit those 40 images, and the reported score is too optimistic.

\vspace{0.2cm}
\begin{graybox}
\raggedright\textbf{\faTrophy\ Your scoreboard} (\texttt{my\_results.csv}) collects every method scored on the
same 40 images: handcrafted features, a CNN from scratch, linear probes, fine-tuning, and later
CLIP and vision-language models.
\end{graybox}
\end{column}
\end{columns}
\end{frame}

% ─────────────────────────────────────────────────────────────────────────────
\begin{frame}{How Many Labels Do You Need?}
\begin{columns}[T]
\begin{column}{0.46\textwidth}
\centering
\begin{tikzpicture}[x=1.0cm, y=4.2cm]
\draw[-{Stealth[length=1.6mm]}, midgray] (0,0.5) -- (5.3,0.5);
\node[font=\scriptsize, color=bodytext] at (2.6,0.39) {labeled images (log scale)};
\draw[-{Stealth[length=1.6mm]}, midgray] (0,0.5) -- (0,1.03) node[above, font=\scriptsize] {balanced accuracy};
\foreach \x/\l in {0.5/10, 1.5/20, 2.5/40, 3.5/80, 4.5/160} \node[below, font=\tiny, color=midgray] at (\x,0.5) {\l};
\fill[wurgreen, opacity=0.18] plot[smooth] coordinates {(0.5,0.52) (1.5,0.6) (2.5,0.72) (3.5,0.8) (4.5,0.84)}
  -- plot[smooth] coordinates {(4.5,0.9) (3.5,0.89) (2.5,0.86) (1.5,0.8) (0.5,0.72)} -- cycle;
\draw[very thick, wurgreen] plot[smooth] coordinates {(0.5,0.62) (1.5,0.7) (2.5,0.79) (3.5,0.845) (4.5,0.87)};
\draw[dotted, midgray] (0,0.5) -- (5,0.5);
\end{tikzpicture}\\
\muted{\footnotesize schematic: mean and range over repeated random draws}
\end{column}
\begin{column}{0.51\textwidth}
A \keyterm{learning curve}: train on random subsets of increasing size, several times per size,
and score each on your codes.

\begin{itemize}
\item Good embeddings often get close to their plateau with a few dozen labels
\item The band is wide at the small end: with 20 labels, \emph{which} 20 you labeled matters a lot
\item With 40 test images, every curve is noisy too
\end{itemize}
\muted{Use learning curves to decide whether more coding is worth it.}
\end{column}
\end{columns}
\end{frame}

% ─────────────────────────────────────────────────────────────────────────────
\begin{frame}{Search by Example, Find Near Duplicates}
\begin{columns}[T]
\begin{column}{0.42\textwidth}
\textbf{Search by example}\\
Compute the cosine similarity between one image and all others, and look at the closest ones.
No labels needed: a quick way to find more examples of a rare kind of image for coding.

\vspace{0.2cm}
\textbf{Near duplicates}\\
Pairs with very high similarity. Embeddings are trained to recognize the same scene under
different crops, so they catch copies that perceptual hashes (module 02) miss.
\end{column}
\begin{column}{0.55\textwidth}
{\footnotesize
\begin{tabular}{@{}>{\raggedright\arraybackslash}p{3.3cm}cc@{}}
\toprule
copy & \shortstack{perceptual\\hash} & embedding \\
\midrule
resized, recompressed & \textcolor{wurgreen}{\faCheck} & \textcolor{wurgreen}{\faCheck} \\
brightened & \textcolor{wurgreen}{\faCheck} & \textcolor{wurgreen}{\faCheck} \\
cropped & \textcolor{midgray}{only lightly} & \textcolor{wurgreen}{\faCheck} \\
mirrored & \textcolor{warnborder}{\faTimes} & \textcolor{wurgreen}{\faCheck} \\
another photo, same moment & \textcolor{warnborder}{\faTimes} & \textcolor{midgray}{close} \\
\bottomrule
\end{tabular}}

\vspace{0.2cm}
\begin{warnbox}
Choose a threshold by \textbf{looking at pairs around it}, and decide what counts as a duplicate
for your question.
\end{warnbox}
\end{column}
\end{columns}
\end{frame}

% ─────────────────────────────────────────────────────────────────────────────
\begin{frame}{Visual Themes: Embed, Reduce, Cluster, Look}
\vspace{0.1cm}
\centering
\begin{tikzpicture}[
  s/.style={draw=wurgreen, fill=lightgreen, rounded corners=3pt, font=\footnotesize, align=center,
            inner sep=3pt, text width=2.05cm, minimum height=1.25cm},
  h/.style={draw=darknavy, fill=darknavy, text=white, rounded corners=3pt, font=\footnotesize,
            align=center, inner sep=3pt, text width=2.05cm, minimum height=1.25cm},
  arr/.style={-{Stealth[length=2mm]}, thick, color=midgray}]
\node[s] (a) {\textbf{embed}\\ DINOv2, 384 numbers};
\node[s, right=0.38cm of a] (b) {\textbf{reduce}\\ UMAP to 5 dimensions};
\node[s, right=0.38cm of b] (c) {\textbf{cluster}\\ HDBSCAN; outliers stay out};
\node[h, right=0.38cm of c] (d) {\textbf{look}\\ the images closest to each center};
\node[h, right=0.38cm of d] (e) {\textbf{name}\\ the themes yourself};
\draw[arr] (a) -- (b); \draw[arr] (b) -- (c); \draw[arr] (c) -- (d); \draw[arr] (d) -- (e);
\end{tikzpicture}

\vspace{0.3cm}
\raggedright
\begin{columns}[T]
\begin{column}{0.48\textwidth}
\begin{graybox}
\textbf{Results depend on settings}, as with topic models. Rerun with other cluster sizes and
seeds; compare runs with the \keyterm{adjusted Rand index}. Trust the themes that survive.
\end{graybox}
\end{column}
\begin{column}{0.48\textwidth}
\begin{warnbox}
\textbf{A cluster is not a frame.} Clusters group what looks alike to the model: product type,
brand, background. Frames come from theory \muted{(Torres, 2024)}. Use clusters to discover, a
codebook to measure.
\end{warnbox}
\end{column}
\end{columns}
\end{frame}

% ─────────────────────────────────────────────────────────────────────────────
\notebookslide{
\begin{itemize}
\setlength{\itemsep}{5pt}
\nbitem{\faTags}{ImageNet labels off the shelf, for protest photos and packaging}
\nbitem{\faCubes}{Embeddings from ResNet-50 and DINOv2}
\nbitem{\faTrophy}{Linear probes, a tuned probe and a fine-tuned ResNet-50, on your scoreboard}
\nbitem{\faChartLine}{A learning curve: how many labels do you need?}
\nbitem{\faSearch}{Search by example and near duplicates}
\nbitem{\faShapes}{Visual themes in cereal packaging, and how stable they are}
\end{itemize}

\vspace{0.35cm}
\muted{\footnotesize part2\_vectors/06\_transfer\_learning\_and\_embeddings/notebooks/\\
06\_transfer\_learning\_and\_embeddings\_exercises.ipynb · 6 exercises · GPU recommended}
}

\end{document}
